\documentclass[10pt,a4paper]{amsart}
\usepackage[margin=19mm]{geometry}
\usepackage{amsmath,amssymb,mathtools}
\usepackage[scr=rsfs,cal=euler]{mathalfa}
\usepackage{xfrac}
\usepackage[unicode,hidelinks,bookmarks=false]{hyperref}
\newcommand{\ie}{i.e.\ }
\newcommand{\cf}{cf.\ }
\newcommand{\resp}{respectively\ }
\newcommand{\Subsection}{\subsection}
\providecommand{\nobreakdash}{\nobreak\hbox{-}}
\setlength{\emergencystretch}{2em}
\multlinegap0pt
\usepackage[shortlabels]{enumitem}
\usepackage{amssymb}
\usepackage{mathtools}
\usepackage{dsfont}
\usepackage[all]{xy}
\newtheorem{lemma}{Lemma}
\newtheorem{corollary}{Corollary}
\newtheorem{proposition}{Proposition}
\def\I{{\mathcal I}}
\def\M{{\mathcal M}}
\DeclareMathOperator{\depth}{depth}
\DeclareMathOperator{\h}{H}
\newcommand{\m}{\mathfrak m}
\title[Author proof 107880]{Multiplicity versus Buchsbaumness of the special fiber ring}
\author{Anoot Kumar Yadav, Kumari Saloni}
\date{Source: arXiv:2102.04218v4 (CC BY 4.0)}
\def\F{{\mathcal F}}
\def\L{{\mathcal L}}
\def\N{{\mathcal N}}
\def\R{{\mathcal R}}
\def\U{{\mathcal U}}
\DeclareMathOperator{\ass}{Ass}
\DeclareMathOperator{\coker}{Coker}
\DeclareMathOperator{\height}{height}
\DeclareMathOperator{\im}{Im}
\newcommand{\p}{\mathfrak p}
\newtheorem{remark}{Remark}
\newtheorem{theorem}{Theorem}
\allowdisplaybreaks
\let\originalcite\cite
\renewcommand{\cite}[2][]{\textup{[author ref.\ \texttt{\detokenize{#2}}\if\relax\detokenize{#1}\relax\else, #1\fi]}}
\makeatletter
\newlabel{the-theorem}{{\textup{source line 937}}{}}
\makeatother
\begin{document}

\maketitle
\noindent\textbf{Reviewed difficulty: H2. Legacy noncounting; count 0.} The H2 burden includes the joint author chain: polynomial fiber realization, the three dedicated kernels and Hilbert-defect identity at 633--724, pure-dimensionality and multiplicity forcing kernel vanishing at 799--868, finite-length and superficial reductions at 576--607, and the dimension/filtration double induction at 905--935. These specialized steps are displayed and remain part of the selected unit\textquotesingle s burden alongside the mature commutative-algebra inputs.\par\medskip

\noindent\emph{Source and licence.} Multiplicity versus Buchsbaumness of the special fiber ring. Version arXiv:2102.04218v4 (CC BY 4.0), distributed by arXiv under the Creative Commons Attribution 4.0 International licence, \url{http://creativecommons.org/licenses/by/4.0/}, as recorded in the arXiv licence metadata for this exact version and shown on the abstract page \url{https://arxiv.org/abs/2102.04218v4}. Full licence notice: \path{licenses/arxiv-2102.04218-CC-BY-4.0.txt}.\par\smallskip


\noindent\textit{Editorial scope.} Selected author statement at source0.tex 901--904 and complete author proof at 905--935. Necessary author inputs and any dedicated auxiliary chain are displayed below. Original mathematical text, hypotheses, formulas and source version are retained.\par\smallskip
\section*{Author context (source 241--242)}
	 \begin{equation}\label{very-1st-eqn(a)}
	 	e_1(\M)-e_1^J(M)\geq 2e_0(\M)-2\ell(M/M_1)-\ell(M_1/(M_2+JM)) \end{equation}

\section*{Author context (source 252--289)}
\section{Notation}\label{notation}

   % \end{equation}
    
    Throughout this paper, let $\I=\{I_n\}_{n\geq 0}$ be a multiplicative $I$-good filtration of $A$ and $\M=\{M_n\}_{n\geq 0}$ be an $I$-good filtration of $M$ as defined in the previous section. The associated graded module $G(\M)=\mathop{\oplus}\limits_{n\geq 0} M_n/M_{n+1}$ and the Rees module $R(\M)=\mathop{\oplus}\limits_{n\geq 0}^\infty M_nt^n$ are modules over the Rees algebra $R(I)=\mathop{\oplus}\limits_{n\geq 0} I^nt^n$. An ideal $J\subseteq I$ is called an {\em $\M$-reduction of $I$} if $M_{n+1}=JM_{n}$ for $n\gg 0.$ We say that $J$ is a {\em minimal $\M$-reduction of $I$} if it is minimal with respect to inclusion. For a reduction $J$ of $I$, let $$\M_J=\{J^nM\}_{n\geq 0},\text{ an $I$-good filtration of $M$ and }$$  $$R(\M_J)= \mathop{\oplus}\limits_{n\geq 0}^\infty J^nM, \text{ a module over } R(J).$$ 
We regard $R(\M)$ and $R(\M_J)$ as submodules of $M[t]$ over polynomial ring $A[t].$ Let $R(\M)_{\geq m}=\mathop{\oplus}\limits_{n\geq m}M_nt^n$ be the graded submodule of $R(\M).$ The unique homogeneous maximal ideals of $R(I)$ and $R(J)$ are denoted by $\U$ and $\xi$ respectively, i.e., $\U=\m R(I)+R(I)_+$ and $\xi = \m R(J)+R(J)_+.$ We write $$F_I(J)=R(J)/IR(J)=\mathop{\oplus}_{n\geq 0} J^n/IJ^n,$$ $$F_{I_1}(\M_J)=R(\M_J)/I_1R(\M_J)=\mathop{\oplus}\limits_{n\geq 0}J^nM/I_1J^nM \text{ and }$$ $$F_\m(\M_J)=R(\M_J)/\m R(\M_J)=\mathop{\oplus}\limits_{n\geq 0}J^nM/\m J^nM.$$


Let $\overline{A}=A/\h_{\m}^0(A)$, $W=\h_{\m}^0(M)$, and $\overline{M}=M/W$ which is an $\overline{A}$-module. Further, $\overline{\M}=\{ \overline{M_n}=(M_n+W)/W\}_{n\geq 0}$ and $\overline{\M_J}=\{(J^nM+W)/W\}_{n\geq 0}$  are $I\overline{A}$-good filtration of $\overline{M}$.
For  $a\in A$, we write $A^\prime=A/(a)$, $M^\prime=M/aM$, $J^\prime=JA^\prime$, $I^\prime=IA^\prime$ and  $\M^\prime=\{M_n^\prime=(M_n+aM)/aM\}_{n\geq 0}$ which is an $I^\prime$-good filtration of $M^\prime$.  

Let $J=(a_1,\ldots,a_d)$ be a minimal $\M$-reduction of $\I$. 
It follows that $\ell(M/JM)<\infty.$ We consider the following conditions on $a_1,\ldots,a_d$ for $M$: 
	\begin{itemize}
	\item[$(C_0)$] The sequence $a_1,\ldots,a_d$ is a $d$-sequence of $M$ as defined in \cite{adt};
	\item[$(C_1)$] The above sequence is an unconditioned strong d-sequence of $M$;
	\item[$(C_2)$]  $(a_1,\ldots, \check{a_i},\ldots,a_d)M:_Ma_i\subseteq IM$   for $1\leq i\leq d$.
	\item[$(C_3)$] $\depth  M>0$.
\end{itemize}
The above conditions were introduced in \cite{ozeki13} when $M=A$. We note that
all the four conditions $(C_0)$-$(C_3)$ are satisfied if $M$ is Cohen-Macaulay. The condition $(C_1)$ is equivalent to saying that $M$ is generalized Cohen-Macaulay and $J$ is a standard parameter ideal for $M$. The role played by regular sequences in the study of Cohen-Macaulay modules is broadly replaced
by $d$-sequences in the case of generalized Cohen-Macaulay modules. 
A sequence $a_1,\ldots,a_s$ of elements of $A$ is said to be a {\textit{d-sequence}} of $M$ if the equality 
\[q_{i-1}M:_M a_ia_j=q_{i-1}M:_M a_j\] holds for $1\leq i\leq j\leq s$ where 
$q_{i-1}=(a_1,\ldots,a_{i-1})$ and $q_0=(0)$. It is said to be an {\textit{unconditioned d-sequence}} if it is a $d$-sequence in any order. Moreover, we say that $a_1,\ldots,a_s$ is an {\textit{unconditioned strong $d$-sequence}} (u.s.d-sequence) if $a_1^{n_1},\ldots,a_s^{n_s}$ is an  unconditioned $d$-sequence for  all integers $n_1,\ldots,n_s>0$. 
Recall that a module is Buchsbaum if and only if every system of parameters of $M$ is u.s.d-sequence of $M$. So,  $(C_0)$ and $(C_1)$ are satisfied by every system of parameters of $M$ if $M$ is a Buchsbaum module. 


\begin{remark}\label{remark-nota}
Suppose  $(C_1)$ is satisfied by $(a_1,\ldots,a_d)$ for $M$. Then $a_i W\subseteq JW\subseteq JM\cap W=0$,  see \cite[Corollary 2.3]{trung-gcm}. This gives 
$ W\subseteq (a_1,\ldots, \check{a_i},\ldots,a_d)M:_Ma_i$ for $1\leq i\leq d.$
\end{remark}

By $\mu(M)$, we denote the minimal number of generators of $M$ and $e_i(\M)$, $1\leq i\leq d$ are the Hilbert coefficients of $\M$, i.e., the unique integers such that $\ell(M/M_n)=\sum\limits_{i=0}^d(-1)^ie_i(\M)\binom{n+d-i-1}{d-i}$ for $n\gg 0.$ When $\M$ is the $I$-adic filtration $\{I^nM\}_{n\geq 0}$, we write $e_i^I(M)$ in place of $e_i(\M)$. 

We assume that the residue field $A/\m$ is infinite. For convenience, we highlight that for most of our results, as in \cite{rossi-valla}, $\M$ is either $$\M_\I=\{I_nM\}_{n\geq 0}\text{ or }$$ $$\M_L=\{M_0=M, ~M_{n+1}=I^nL ~\text{ for } ~n\geq 0\} \text{ where } L=\m M 
\text{ and }$$
through out the paper, $J=(a_1,\ldots,a_d)\subseteq I$ denote either a minimal $\M_\I$-reduction of $I$ or a minimal $\M_L$-reduction of $I$ which will be clear from the context. 

\section*{Author context (source 325--334)}
\section{Preliminaries}\label{Section-Pre}
In this section, we present some auxiliary results that will be used in subsequent sections. We discuss most of the proofs for the reader's convenience and mention the references for $I$-adic cases of some results.
\subsection{Sally module of filtration}
To study blow-up algebras, Vasconcelos \cite{Vas94} introduced the notion of Sally modules in the ring case. 
Further, Rossi-Valla in \cite[Chapter 6]{rossi-valla}, extended the definition for an $I$-good filtration of a module which we recall below. The Sally module, $S_J(\M)$ corresponding to the filtration $\M_\I$ or $\M_L$ can be defined as the cokernel of the following maps of graded $R(J)$-modules respectively:
\begin{align}\label{eq_1}
 &0\longrightarrow I_1R(\M_J)\longrightarrow \R(\M)_{+}(+1)\longrightarrow S_J(\M)\longrightarrow 0 \\
 &0\longrightarrow \m R(\M_J)\longrightarrow \R(\M)_{+}(+1)\longrightarrow S_J(\M)\longrightarrow 0
\end{align}
where $J=(a_1,\ldots,a_d)\subseteq I$ is a minimal $\M_\I$-reduction or $\M_L$-reduction  of $I$ respectively. Therefore, for both the filtration, Sally module is a finitely generated graded $R(J)$-module and $\dim_{R(J)} S_J(\M)\leq d$.  

\section*{Author context (source 352--437)}
\begin{lemma}\label{lemma-3.1(a)} Let $M$ be a finitely generated $A$-module of dimension $d$.  Suppose $\M=\M_\I$ and $J\subsetneq I_1$ or $\M=\M_L$ and $J\subsetneq \m.$ 
Then there exists an exact sequence of graded $R(J)$-modules
$$
    R(\M_J)(-1)^\mu\xrightarrow{\phi}R(\mathcal{M})/ R(\M_J)\to S_J(\M)(-1)\to 0
    $$ 
    where  $\mu=\mu(\m \setminus J)$ when $\M=\M_L$ and  $\mu=\mu(I_1\setminus J)$ when $\M=\M_\I.$ 
\end{lemma}
\begin{proof}
 Suppose $\M=\M_L$ and  $\m =J+(y_1,\ldots,y_\mu)$, where $y_1,\ldots,y_\mu \in \m$. Consider the graded map 
	$$R(\M_J)(-1)^\mu\xrightarrow{\phi} R(\mathcal{M})/R(\M_J)$$
	defined as $\phi(\alpha_1,\ldots,\alpha_\mu )=\overline{\mathop\sum\limits_{i=1}^\mu y_i \alpha_i }\in R(\mathcal{M})/R(\M_J)$, where $\alpha_i \in J^{n-1}M $ and $y_1,\ldots,y_{\mu}\in \m $. Then
	$[\coker \phi]_{n}=\m I^{n-1}M/(J^nM+J^{n-1}\m M)=\m I^{n-1}M/J^{n-1}\m M=(S_J(\M)(-1))_{n}$ for $n\geq 1$ and
	$[\coker \phi]_{0}=0$. Therefore,
	$\coker \phi\cong S_J(\M)(-1)$ as graded $R(J)$-module.

    Now assume $\M=\M_\I$ and $I_1=J+(y_1,\ldots,y_\mu)$, where $y_1,\ldots,y_\mu \in I_1$. As earlier, define the graded map $R(\M_J)(-1)^\mu\xrightarrow{\phi} R(\mathcal{M})/R(\M_J)$ by  $\phi(\alpha_1,\ldots,\alpha_\mu )=\overline{\mathop\sum\limits_{i=1}^\mu y_i \alpha_i}$. Then $[\coker \phi]_{n}=I_{n}M/(J^nM+J^{n-1}I_1M)=I_{n}M/J^{n-1}I_1M=(S_J(\M)(-1))_{n}$ for $n\geq 1$ and
	$[\coker \phi]_{0}=0$ which gives 
	$\coker \phi\cong S_J(\M)(-1)$ as graded $R(J)$-module.
\end{proof}
% \begin{lemma}\label{lemma-3.1(b)} Let $(A,\m)$ be a Noetherian local ring and $M$ be a finitely generated $A$-module. Let $\M=\{I_nM\}$ be an $I$-good filtration of $M$ induced by an $I$-good filtration $\I=\{I_n\}$ of $A.$
%  Assume that $J \subsetneq I_1$ and $\mu=\mu(I_1\setminus J)$. Then there exists an exact sequence
%     $$
%     R(\M_J)(-1)^\mu\xrightarrow{\phi}R(\mathcal{M})\setminus R(\M_J)\to S_J(\M)(-1)\to 0
%     $$
%     of graded $R(J)$-modules.

% \end{lemma}
%%\begin{proof}
% Let $I_1=J+(x_1,\ldots,x_\mu)$, where $x_1,\ldots,x_\mu \in I_1$. Consider the map 
%	$$R(\M_J)(-1)^\mu\xrightarrow{\phi} R(\mathcal{M})/R(\M_J)$$
%	defined as $\phi(\alpha_1,\ldots,\alpha_\mu )=\overline{\mathop\sum\limits_{i=1}^\mu \alpha_ix_i t}\in R(\mathcal{M})/R(\M_J)$ where $\alpha_i \in J^{n-1} M$. Then
%	$[\coker \phi]_{n}=I_{n}M/(J^nM+J^{n-1}I_1M)=I_{n}M/J^{n-1}I_1M=(S_J(\M)(-1))_{n}$ for $n\geq 1$ and
%	$[\coker \phi]_{0}=0$. Therefore,
%	$\coker \phi\cong S_J(\M)(-1)$ as graded $R(J)$-module.
%\end{proof}
% \begin{lemma}\label{lemma-2.2}
% 	Suppose that the conditions $(C_0)$ and $(C_2)$ are satisfied for an $I$-good filtration $\M$ and $A$. Suppose $I_1$ is an ideal of $A$. Then $R(\N)/I_1R(\N)=\F^\prime\cong (M/I_1M)[X_1,\ldots,X_d]$ and $R(J)/IR(J)=\F \cong A/I[X_1,\ldots,X_d]$ as graded $A$-algebras, respectively. 
% 	In particular, $\F$ and $\F^\prime$ are Cohen-Macaulay rings of dimension $d$.
% \end{lemma}
% \begin{proof}First we prove for $\M=\{I_nM\}$. In \cite[Theorem 3.3]{trung-abs}, Trung show that $R(\N)=Sym_A(JM)$, where $Sym_A(JM)$ denotes the symmetric algebra of $J$ with respect to $M$. Notice that 
% $$
% \frac{R(\N)}{I_1R(\N)}= \frac{A}{I_1} \otimes Sym_A(JM)\cong Sym_{A/I_1}\Big(\frac{A}{I_1} \otimes JM\Big)=Sym_{A/I_1}\Big(\frac{JM}{I_1JM}\Big)=Sym_{A/I_1}\Big(\frac{JM}{JI_1M}\Big)  
% $$
% as graded $A/I_1$-algebras and we have 
% $$
% R(\N)/I_1R(\N)=M/I_1M[X_1,\ldots, X_d]
% $$
% \textcolor{red}{because by condition $(C_2)$ the $A/I_1$-module $JM /I_1JM$ is free of rank $d$.}\\
% Now we prove for $\M=\M_L$, where $L=\m M$. In \cite[Theorem 3.3]{trung-abs}, Trung show that $R(\N)=Sym_A(JM)$, where $Sym_A(JM)$ denotes the symmetric algebra of $J$ with respect to $M$. Notice that 
% $$
% \frac{R(\N)}{\m R(\N)}= \frac{A}{\m} \otimes Sym_A(JM)\cong Sym_{A/\m}\Big(\frac{A}{\m} \otimes JM\Big)=Sym_{A/\m}\Big(\frac{JM}{\m JM}\Big)=Sym_{A/\m}\Big(\frac{JM}{J\m M}\Big) 
% $$
% as graded $A/\m$-algebras and we have 
% $$
% R(\N)/\m R(\N)=M/\m M[X_1,\ldots, X_d]
% $$
% \textcolor{red}{because by condition $(C_2)$ the $A/\m$-module $JM /\m JM$ is free of rank $d$.}\\
	 
% \end{proof}
\begin{lemma}\label{lemma-2.2(a)}
	Let $M$ be a finitely generated flat $A$-module of dimension $d>0$. Suppose $(C_0)$ and $(C_2)$ conditions are satisfied by the sequence $a_1,\ldots,a_d$ for both $A$ and $M$. Then 
    \begin{enumerate}
        \item $R(\M_J)/I_1R(\M_J)=\F_{I_1}(\M_J)\cong (M/I_1M)[X_1,\ldots,X_d]$ as graded $R(J)$-module,
        \item $R(J)/IR(J)=\F_I(J) \cong A/I[X_1,\ldots,X_d]$ as graded $A$-algebra and 
        \item $R(\M_J)/\m R(\M_J)=\F_{\m}(\M_J)\cong (M/\m M)[X_1,\ldots,X_d]$ as graded $R(J)$-module. 	
    \end{enumerate}  
	In particular, $\F_I(J)$ is Cohen-Macaulay ring of dimension $d$, $\F_{I_1}(\M_J)$ and  $\F_\m(\M_J)$ are Cohen-Macaulay $R(J)$-modules of dimension $d$.
\end{lemma}
\begin{proof}
    Since $M$ is a flat $A$-module, we have $R(\M_J)\cong R(J) \otimes_A M.$ By condition $(C_0)$, we get that $R(J)\cong Sym_A(J)$, where $Sym_A(J)$ denotes the symmetric algebra of $J$. Notice that 
\begin{eqnarray*}
\frac{R(\M_J)}{I_1R(\M_J)}\cong \frac{A}{I_1}\otimes_A R(J)\otimes_A M &\cong&  \frac{A}{I_1} \otimes_A Sym_A(J)\otimes_A M\\
&\cong& Sym_{A/I_1}\Big(\frac{A}{I_1} \otimes_A J\Big)\otimes_A M\\
&=&Sym_{A/I_1}\Big(\frac{J}{I_1J}\Big)\otimes_A M \text{ as graded $A/I_1$-module and } 
\end{eqnarray*}
\begin{eqnarray*}
    \frac{R(\M_J)}{\m R(\M_J)}\cong \frac{A}{\m}\otimes_A R(J)\otimes_A M &\cong&  \frac{A}{\m} \otimes Sym_A(J)\otimes_A M\\
    &\cong& Sym_{A/\m}\Big(\frac{A}{\m} \otimes_A J\Big)\otimes_A M\\&=&Sym_{A/\m}\Big(\frac{J}{\m J}\Big)\otimes_A M  \text{ as graded $A/\m$-module. }
\end{eqnarray*}

By condition $(C_2)$, $J/I_1J$ and $J/\m J$ are free  of rank $d$ over $A/I_1$ and $A/\m$ respectively. Hence $Sym_{A/I_1}(J/I_1J)\cong A/I_1[X_1,\ldots,X_d]$ and $Sym_{A/\m}(J/\m J)\cong A/\m[X_1,\ldots,X_d]$.   Therefore, we have 
\begin{align*}
R(\M_J)/I_1R(\M_J)&\cong A/I_1[x_1,\ldots,X_d]\otimes_A M \cong M/I_1M[X_1,\ldots, X_d] \text{ and }\\
R(\M_J)/\m R(\M_J)&\cong A/\m [x_1,\ldots,X_d]\otimes M \cong M/\m M[X_1,\ldots, X_d]\qedhere
\end{align*}
\end{proof}

\section*{Author context (source 473--622)}
\begingroup\footnotesize
\begin{lemma}\cite[Corollary 6.2]{rossi-valla}\label{lemma-2.3(a)}
    Let $M$ be a Cohen-Macaulay $A$-module of dimension $d$ and $L$ be a submodule of $M$. Then $S_J(\M_L)=0$ or $\dim S_J(\M_L)=d.$ 
    
    
   %% Let $\M=\M_L$ and $J\subseteq I$ be an ideal generated by a maximal $\M$-superficial sequence for $I$, then $S_J(\M)=0$ or $\dim S_J(\M)=d.$ 
    % \begin{enumerate}
    %     \item $\dim S_J(\M)=d$ provided it is not the trivial module.
    %     \item $e_0(S_J(\M))=e_1(\M)-\ell(L/JM)$ and $e_i(S_J(\M))=e_{i+1}(\M)$ for every $i>0$
    % \end{enumerate}
\end{lemma}
To determine the dimension of $S_J(\M)$ when $\M=\M_{\I}$, we state the following lemma. 
 \begin{lemma}\label{lemma-2.3(b)} Suppose $\depth A>0$, $M$ is a finitely generated faithful $A$-module of dimension $d$ and $\depth M>0.$ Suppose that the conditions $(C_0)$ and $(C_2)$  are satisfied by $a_1,\ldots,a_d$ for both $A$ and $M$.  If $R(\M_J)$ satisfies Serre's condition $(S_2)$, then $\ass_{R(J)}(S_J(\M_\I))\subseteq \{\m R(J)\}$. Consequently, $\dim_{R(J)}(S_J(\M_\I))=d$ if $S_J(\M_\I)\neq 0$.
 \end{lemma}
 % \begin{lemma} 
 % ??Let $(A,\m)$ be a Noetherian local ring, $J=(a_1,\ldots,a_d) $ be a minimal reduction of $\I$ and $M$ be a finitely generated faithful flat $A$-module. Suppose $(C_0)$ and $(C_2)$ condition are satisfied by $a_1,\ldots,a_d$ for both $A$ and $M$. 
 % 	Also, $(C_3)$ are satisfied by $M$ and $A$. Let $\I=\{I_n\}$ be an $I$-good filtration of $A$ and $\M=\{I_nM\}$ be an $I$-good filtration of $M$. If $R(\M_J)$ satisfies Serre's condition $(S_2)$. Then $\ass_{R(J)}(S_J(\M))\subseteq \{\m R(J)\}$.
 % 	Hence, $\dim_{R(J)}(S_J(\M))=d$ if $S_J(\M)\neq 0$.
 % \end{lemma}
 % For $\I=\{I^n\}$ and $\I=\{\overline{I^n}\}$, the above result is discussed in \cite[Lemma 2.3]{goto-ozeki-ncm-2010} and  \cite[Lemma 2.2]{phuong-normal-2018} respectively. 
 %  We provide a proof for an $I$-good filtration.
 \begin{proof}
 	Let $P\in\ass_{R(J)}(S_J(\M_\I))$. Since $\m^l S_J(\M_\I)=0$ for $l\gg 0$, we have $\p=\m R(J)\subseteq P.$ Suppose $\p\neq P$. 
 	Note that $\depth_{R(J)_P} (R(\M_\I)_{+}(+1))_P\geq 1$
 	since $a_1\in \p\subseteq P$
 	is a non-zero-divisor on $R(\M_\I)_{+}(+1)=\mathop\oplus\limits_{j\geq 1}I_jMt^{j-1}$. Applying the depth Lemma on the exact sequence
 	$$0\longrightarrow (I_1 R(\M_J))_P\longrightarrow (R(\M_\I)_{+}(+1))_P\longrightarrow (S_J(\M_\I))_P\longrightarrow 0,$$
 	of graded $R(J)_P$-modules. We get $$\depth_{R(J)_P} (I_1 R(\M_J))_P\geq \min\{ \depth_{R(J)_P} (R(\M_\I)_{+}(+1))_P,\depth_{R(J)_P} (S_J(\M_\I))_P+1\}\geq 1$$ and 
 	$\depth_{R(J)_P} (S_J(\M_\I))_P=0\geq \min\{\depth_{R(J)_P} (R(\M_\I)_{+}(+1))_P,\depth_{R(J)_P} (I_1 R(\M_J))_P-1\}=\depth_{R(J)_P} (I_1 R(\M_J))_P-1$. Hence $\depth_{R(J)_P} (I_1 R(\M_J))_P=1$.
 	Now consider the exact sequence
 	$$0\longrightarrow (I_1 R(\M_J))_P\longrightarrow R(\M_J)_P\longrightarrow (R(\M_J)/I_1 R(\M_J))_P\longrightarrow 0.$$
 	Since $\p\neq P$, $\height~P\geq 2$.
 	So $\dim R(J)_P\geq 2$. Since $M$ is faithful $A$-module, hence $R(\M_J)$ is faithful $R(J)$-module. Therefore $\dim R(\M_J)_P\geq 2$ which implies that $\depth R(\M_J)_p\geq \min\{2,\dim R(\M_J)_P\}=2$ as $R(\M_J)$ satisfies $(S_2)$. 
 	Again by depth lemma, $1=\depth_{R(J)_P} (I_1 R(\M_J))_P\geq \min\{\depth_{R(J)_P} R(\M_J)_P,\depth_{R(J)_P} (R(\M_J)/I_1 R(\M_J))_P+1 \}=\depth_{R(J)_P} (R(\M_J)/I_1 R(\M_J))_P+1$ which implies
 	$\depth_{R(J)_P} (R(\M_J)/I_1 R(\M_J))_P=0.$ By Lemma \ref{lemma-2.2(a)}, $R(\M_J)/I_1 R(\M_J)$ is Cohen-Macaulay which means $\dim (R(\M_J)/I_1 R(\M_J))_P=\depth_{R(J)_P} (R(\M_J)/I_1 R(\M_J))_P=0.$ Thus 
 	$P\in\min_{R(J)}(R(\M_J)_P/I_1 R(\M_J)_P)=\{\m R(J)\}$ which is a contradiction.
 \end{proof}
The Hilbert coefficients $e_i(S_J(\M))$ of the Sally module $S_J(\M)$ and the coefficients $e_i(\M)$ of $\M$ are related through the following relations. Here $\M=\M_L$ or $\M=\M_\I.$  We have
\[\ell(M/M_{n+1})=\ell(M/J^n M_1)-\ell(S_J(\M)_n)= \ell(M/J^nM)+\ell(J^nM/J^n M_1)-\ell(S_J(\M)_n)\]
for $n\geq 0$. 
%%The next result is a generalization of \cite[Proposition 2.4]{ozeki13}. 
\begin{proposition}\label{proposition:p1} Let $M$ be a finitely generated faithful $A$-module and $\M=\M_L$ or $ \M=\M_\I.$  Then
	\begin{equation*}
		e_1(\M)\geq  \begin{cases} e_1^J(M)+e_0^J(M)-\ell_A(M/M_1)+e_0(S_J(\M)) &\text{ if } \dim S_J(\M)=d.\\
			e_1^J(M)+e_0^J(M)-\ell_A(M/M_1) & \text{ if } \dim S_J(\M)\leq d-1.
		\end{cases}
	\end{equation*}
\end{proposition}
\begin{proof}
	For $n\geq 0$, when $M_1=I_1M$, we have $J^nM/J^n M_1= (R(\M_J)/I_1 R(\M_J))_n$ and $R(\M_J)/I_1 R(\M_J)$ is a homomorphic image of $M/M_1[X_1,\ldots,X_d]$ or when $M_1=\m M$, we have $J^nM/J^n M_1= (R(\M_J)/\m R(\M_J))_n$ and $R(\M_J)/\m R(\M_J)$ is a homomorphic image of $M/M_1[X_1,\ldots,X_d]$. So for all $n\geq 0$, 
	\begin{equation}
		\ell(J^nM/J^{n} M_1)\leq \ell(M/M_1)\binom{n+d-1}{d-1}.\label{equation:step-hilbFunfibercone}
	\end{equation}
	% % % %  the multiplicity of T/\I_1 T is less than \ell(R/\I_1), see Vasconcelos survey paper
	There exists an integer $N$ such that for all $n\geq N$,
	\begin{align*}
		\ell(M/J^nM)&=\mathop\sum\limits_{i=0}^d (-1)^ie_i^J(M)\binom{n-1+d-i}{d-i}\\
		&=e_0^J(M)\binom{n+d}{d}-(e_0^J(M)+e_1^J(M))\binom{n+d-1}{d-1}\\
		&\qquad+\mathop\sum\limits_{i=2}^d(-1)^i(e_{i-1}^J(M)+e_i^J(M))\binom{n+d-i}{d-i}.
	\end{align*}
	Thus, for all $n\geq N$,
	\begin{align*}
		\begin{split}
			\ell(M/M_{n+1})+\ell(S_J(\M)_n)&=\ell(M/J^nM)+\ell(J^nM/J^n M_1)\\
			&\leq e_0^J(M)\binom{n+d}{d}-(e_0^J(M)+e_1^J(M))\binom{n+d-1}{d-1}+\mathop\sum\limits_{i=2}^d(-1)^i\\
			&\qquad\big(e_{i-1}^J(M)+ e_i^J(M)\big)\binom{n+d-i}{d-i}+\ell(M/M_1)\binom{n+d-1}{d-1}\\
			&\leq e_0^J(M)\binom{n+d}{d}-\Big\{ e_0^J(M)+e_1^J(M)-\ell(M/M_1)\Big\}\binom{n+d-1}{d-1}\\
			&\qquad+\mathop\sum\limits_{i=2}^d(-1)^i(e_{i-1}^J(M)+e_i^J(M))\binom{n+d-i}{d-i} 
		\end{split}
	\end{align*}
	Now, the conclusion  follows on comparing the coefficients.
\end{proof}
\begin{proposition}\label{proposition-2.4}    Let $M$ be a finitely generated faithful flat $A$-module and $\M=\M_L$ or $\M=\M_\I.$ Suppose $(C_0)$ and $(C_2)$ conditions are satisfied by $a_1,\ldots,a_d$ for both $A$ and $M$. Then,  for all $n\geq 0,$ 
	\begin{align}\label{newlabel-eq-1}
		\begin{split}
			\ell_A(M/M_{n+1})&= e_0^J(M)\binom{n+d}{d}-\Big\{ e_0^J(M)+e_1^J(M)-\ell(M/M_1)\Big\}\binom{n+d-1}{d-1}\\
			&\qquad+\mathop\sum\limits_{i=2}^d(-1)^i(e_{i-1}^J(M)+e_i^J(M))\binom{n+d-i}{d-i}-\ell_A(S_J(\M)_{n}).
		\end{split}
	\end{align}
\end{proposition}
\begin{proof} Suppose $(C_0)$ and $(C_2)$ condition are satisfied by $a_1,\ldots,a_d$ for both $A$ and $M$. Then by \cite[Proposition 4.1]{trung-abs}, $N=0$ in the proof of Proposition \ref{proposition:p1}. We have $J^nM/J^n M_1= (R(\M_J)/I_1 R(\M_J))_n$ when $\M=\M_\I$ and $J^nM/J^n M_1= (R(\M_J)/\m R(\M_J))_n$ when $\M=\M_L$. By Lemma \ref{lemma-2.2(a)}, it is isomorphic to $M/M_1[X_1,\ldots,X_d]$.  Hence, equality holds in
	\eqref{equation:step-hilbFunfibercone} for all $n\geq 0.$
\end{proof}
In \eqref{newlabel-eq-1}, we may write
$$\ell_A((S_J(\M)_n)=e_0(S_J(\M))\binom{n+s-1}{s-1}-e_1(S_J(\M))\binom{n+s-2}{s-2}+\ldots+(-1)^{s-1}e_{s-1}(S_J(\M))$$
for $n\gg 0$ where $s=\dim S_J(\M)$. Then, comparing the coefficients of both sides, we get the following corollary.
\begin{corollary}\label{corr-2.5}
Let the hypothesis be the same as in Proposition \ref{proposition-2.4}. Let $\dim S_J(\M)=s$. 
	\begin{enumerate}[label=(\roman*)]
		\item Suppose $s=d$. Then 
		\begin{enumerate}
			\item $e_1(\M)=e_0^J(M)+e_1^J(M)-\ell(M/M_1)+e_0(S_J(\M))$ and 
			\item $e_i(\M)=e_{i-1}^J(M)+e_i^J(M)+e_{i-1}(S_J(\M))$ for $2\leq i\leq d$.
		\end{enumerate}
		\item Suppose $s<d$. Then 
		\begin{enumerate}
			\item $e_1(\M)=e_0^J(M)+e_1^J(M)-\ell(M/M_1)$,
			\item $e_i(\M)=e_{i-1}^J(M)+e_i^J(M)$ for $2\leq i\leq d-s$ and
			\item $e_i(\M)=e_{i-1}^J(M)+e_i^J(M)+(-1)^{d-s}e_{i-d+s-1}(S_J(\M))$ for all $d-s+1\leq i\leq d.$
		\end{enumerate}
	\end{enumerate}
\end{corollary}
\subsection{Reduction steps} In the theory of Hilbert polynomials, it is very useful to reduce a result (i) to the case when $\depth M>0$ and (ii) to a lower-dimensional case for applying induction. The following results allow us to do both of these reductions in most of our results.

\begin{lemma}\label{lemma-reduction-steps}
	\begin{enumerate}[label=(\roman*)]
		\item  If $(C_1)$ and $(C_2)$ conditions are satisfied by $a_1,\ldots,a_d$ for $M$, then $(C_1)$ and $(C_2)$ conditions are satisfied by $a_1,\ldots,a_d$ for $\overline{M}$. 
		
		\item \label{lemma-4.1} Suppose that $d\geq 2$. Let $a\in J\setminus \m J$ be a superficial element for $\M$ and $\M_J$.  Then  $2e_0(\M)-e_1(\M)+e_1^J(M)=2\ell_A(M/M_1)+\ell_A(M_1/(M_2+JM))$ if and only if
		$2e_0(\M^\prime)-e_1(\M^\prime)+e_1^J(M^\prime)=2\ell_{A}(M^\prime/{M}^\prime_1)+\ell_{A}({M}^\prime_1/({M}^\prime_2+JM^\prime)).$
		
		\item \label{lemma-4.2} $2e_0(\M)-e_1(\M)+e_1^J(M)=2\ell_A(M/M_1)+\ell_A(M_1/(M_2+JM))$ if and only if  
		$2e_0(\overline{\M})-e_1(\overline{\M})+e_1^J(\overline{M})=2\ell_A(\overline{M}/\overline{M_1} )+\ell_A(\overline{M_1} /(\overline{M_2} +J\overline{M}))$ and $W\subseteq M_2+JM$.
	\end{enumerate}
\end{lemma}
\begin{proof}
	\begin{enumerate}[label=(\roman*)]
		\item
		Suppose $a_1,\ldots,a_d$ satisfies  $(C_1)$ for $M$. Then it follows from \cite[Theorem 2.1]{trung-gcm} and \cite[Lemma 1.6]{trung-gcm} that $\overline{M}$ is generalized Cohen-Macaulay with $J\overline{A}$ a standard parameter ideal in $\overline{M}$ and $JM\cap W=0$ from \cite[Corollary 2.3]{trung-gcm}. Now let  $(a_1,\ldots,\check{a_i},\ldots,a_d)M:_M a_i\subseteq IM$ and $r+W\in (a_1,\ldots,\check{a_i},\ldots,a_d)\overline{M}:_{\overline{M}} a_i$.  Then $ra_i\in (a_1,\ldots,\check{a_i},\ldots,a_d)M+W$. Suppose $ra_i=s+w_0$ with $s\in (a_1,\ldots,\check{a_i},\ldots,a_d)M$ and $w_0\in W$. Then
		$w_0\in JM\cap W=0$ which implies that $ r\in (a_1,\ldots,\check{a_i},\ldots,a_d)M: a_i\subseteq IM.$ Therefore,   $(a_1,\ldots,\check{a_i},\ldots,a_d)\overline{M}:_{\overline{M}} a_i \subseteq I\overline{M}.$
        \vspace{0.25cm}
        
		\item We have $e_i(\M)=e_i(\M^\prime)$, $e_i^J(M)=e_i^J(M^\prime)$ for $0\leq i\leq d-2$, $e_{d-1}(\M)=e_{d-1}(\M^\prime)+(-1)^d\ell(0:_M a)$ and $e_{d-1}^J(M)=e_{d-1}^J(M^\prime)+(-1)^d\ell(0:_M a)$. Therefore, $2e_0(\M)-e_1(\M)+e_1^J(M)= 2e_0(\M^\prime)-e_1(\M^\prime)+e_1^J(M^\prime)$. Further, $2\ell_A(M/M_1)+\ell_A(M_1/(M_2+JM))=2\ell_{A}(M^\prime/M_1^\prime)+\ell_{A}(M_1^\prime/(M_2^\prime+JM^\prime))$.

         \vspace{0.25cm}
		\item 
		We have $e_i(\M)=e_i(\overline{\M})$ and $e_i^J(M)=e_i^J(\overline{M})$ for $0\leq i\leq d-1$, $e_d(\M)=e_d(\overline{\M})+(-1)^d\ell_A(W)$ and $e_d^J(M)=e_d^J(\overline{M})+(-1)^d\ell_A(W)$. Thus,
		\begin{align*}
			2e_0(\M)-e_1(\M)+e_1^J(M)&= 2e_0(\overline{\M})-e_1(\overline{\M})+e_1^J(\overline{M})\\
			&\leq 2\ell_A(\overline{M}/\overline{M_1})+\ell_A(\overline{M_1}/(\overline{M_2}+J\overline{M})) \mbox{ by \eqref {very-1st-eqn(a)}}\\
			&\leq 2\ell_A(M/M_1)+\ell_A(M_1/(M_2+JM)).
		\end{align*}
		Therefore, $2e_0(\M)-e_1(\M)+e_1^J(M)= 2\ell_A(M/M_1)+\ell_A(M_1/(M_2+JM))$ if and only if $2e_0(\overline{\M})-e_1(\overline{\M})+e_1^J(\overline{M})=2\ell_A(\overline{M}/\overline{M_1})+\ell_A(\overline{M_1}/(\overline{M_2}+J\overline{M})),$ $\ell_A(\overline{M_1}/(\overline{M_2}+J\overline{M}))=\ell_A(M_1/(M_2+JM))$  and $\ell_A(\overline{M}/\overline{M_1})=\ell_A(M/M_1)$. 
		Note that the last two equalities hold if and only if $W\subseteq M_2+JM$. \qedhere
	\end{enumerate}
\end{proof}
We end this section with the following remark. 
\begin{remark}\label{remark-I1=Q-case}
	Suppose $e_1(\M)-e_1^J(M)=2(e_0^J(M)-\ell(M/JM)).$ Then the following conditions hold
    :
	\begin{enumerate}[label=(\roman*)]
		\item  $e_1(\M)=e_1^J(M)$,
		\item $M$ is Cohen-Macaulay $A$-module and 
		\item $\M=\M_J.$
	\end{enumerate}
\end{remark}
\begin{proof}
	We have $e_0(\M)-\ell(M/M_1)\leq e_1(\M)-e_1^J(M)=2(e_0^J(M)-\ell(M/JM))$ which gives 
	$e_0^J(M)-2\ell(M/JM)\geq -\ell(M/M_1)\Rightarrow e_0^J(M)-\ell(M/JM)\geq \ell(M_1/JM)\geq 0.$ On the other hand, $e_0^J(M)\leq \ell(M/JM)$. Therefore, we get $e_0^J(M)=\ell(M/JM)$. 
	Thus $M$ is Cohen-Macaulay $A$-module, see \cite[Corollary 4.7.11]{bruns-herzog} and $e_1(\M)=e_1^J(M)=0.$ By \cite[Theorem 2.7]{rossi-valla} $G(\M)$ is Cohen-Macaulay $A$-module and $M_n\subseteq JM$ for all $n.$ Hence by Valabrega-Valla criteria $M_n=M_n\cap JM= JM_{n-1}$ for all $n\geq 1.$ This gives that $\M=\M_J.$
\end{proof}
\par\endgroup

\section*{Author context (source 633--724)}
Now onwards, we restrict ourselves to write $\M=\{M_n\}_{n\geq 0}$ for either the filtration $\M_\I$ or $\M_L$. 
For a minimal $\M$-reduction $J$ of $I$ , we consider the canonical graded $R(J)$-module homomorphisms: $$\F_{I_1}(\M_J)\to G(\M_\I)=R(\M_\I)/\R(\M_\I)_{+}(1) \text{ and }$$  $$\F_{\m}(\M_J)\to G(\M_L)=R(\M_L)/\R(\M_L)_{+}(1)$$ which give  exact sequences 
	\begin{eqnarray} \label{equation:step1}
		0\to K^{(2)}(\M_\I)\to \F_{I_1}(\M_J)\to G(\M_\I)\to R(\M_\I)/(R(\M_\I)_{+}(1)+R(\M_J))\to 0 \text { and }\\
       \label{equation:step1(a)} 0\to K^{(2)}(\M_L)\to \F_{m}(\M_J)\to G(\M_L)\to R(\M_L)/(R(\M_L)_{+}(1)+R(\M_J))\to 0
	\end{eqnarray}    
	of graded $R(J)$-modules where $K^{(2)}(\M_\I)$ and $K^{(2)}(\M_L)$ denote the respective kernels.  Now, let \begin{equation}
    \mu= \begin{cases}
\mu(I_1\setminus J) & \mbox{ if } \M=\M_\I \\  
\mu(\m/J) & \mbox{ if } \M=\M_L 
    \end{cases}
    \end{equation}
and $y_1,\ldots,y_{\mu}$, elements of $A$, such that $I_1=J+(y_1,\ldots,y_\mu)$ when $\M=\M_\I$ and $\m=J+(y_1,\ldots,y_\mu)$ when $\M=\M_L$.  Consider the graded homomorphism of graded $R(J)$-modules $$((R(\M_J)/IR(\M_J))(-1))^\mu\xrightarrow{\phi} R(\M)/(R(\M)_{+}(1)+R(\M_J))$$ defined as
	$$\phi(\alpha_1,\ldots,\alpha_\mu )=\overline{\mathop\sum_{i=1}^\mu y_i\alpha_i  }\in R(\M)/(R(\M)_{+}(1)+R(\M_J))$$
    where $\alpha_i\in R(\M_J)/I R(\M_J)$ for $1\leq i\leq \mu.$ Then 
$[\coker \phi]_n=\frac{M_n/(M_{n+1}+J^nM)}{( J^{n-1}M_1+M_{n+1})/(M_{n+1}+J^nM)}\cong M_n/(M_{n+1}+ J^{n-1}M_1)=(S_J(\M)/\L(\M))_{n-1}$ where $\L(\M)=\mathop\oplus\limits_{n\geq 1}\frac{M_{n+2}+J^{n}M_1}{ J^{n}M_1}$.  Thus 
	$\coker\phi\cong (S_J(\M)/\L(\M))(-1)$ and we get the exact sequence
	\begin{equation}\label{equation:step2}
    \begin{aligned}
        0\to \ker\phi&\to ((R(\M_J)/IR(\M_J)) (-1))^\mu\xrightarrow{\phi} R(\M)/(R(\M)_+(1)+R(\M_J))\\
        &\to (S_J(\M)/\L(\M))(-1)\to 0.
    \end{aligned}
	\end{equation}
Now we state the main theorem of this section. For the proof, we borrow the ideas from \cite{ozeki13}.  
 \begin{theorem}\label{lemma-3.3} 
	 Let $M$ be a finitely generated faithful flat $A$-module, $\M=\{M_n\}_{n\geq 0}$ be one of  the filtration $\M_\I$ or $\M_L$  and $J=(a_1,\ldots,a_d)$ be a minimal $\M$-reduction of $I$. We assume $J\subsetneq I_1$ in case of $\M=\M_\I$ and $J\subsetneq \m$  in case of $\M=\M_L$. Suppose the conditions $(C_0)$ and $(C_2)$ are satisfied by $a_1,\ldots,a_d$ for both $A$ and $M$. Then 
    \begin{align*}		\ell_A(M_n/M_{n+1})&=\Big \{\ell_A(M/M_1)+\ell_A(M_1/(M_2+JM))\Big\}\binom{n+d-1}{d-1}
		-\ell(M_1/({M_2}+JM))\\&\qquad \binom{n+d-2}{d-2}
		 +\ell_A(S_J(\M)_{n-1})-\ell_A([K^{(1)}(\M)]_{n})-\ell_A([K^{(2)}(\M)]_{n})-\\&\qquad\ell_A([K^{(3)}(\M)]_{n-1})
	\end{align*}
	for all $n\geq 0$ and for some   graded $R(J)$-modules $K^{(i)}(\M)$, $1\leq i\leq 3$. Furthermore, 
	$$
		e_1(\M)-e_1^J(M)= 2 e_0^J(M)-  2\ell_A(M/M_1)-\ell_A(M_1/(M_2+JM))+\mathop\sum\limits_{i\in\Gamma}e_0(K^{(i)}(\M))
	$$
    where $\Gamma=\{i\mid 1\leq i\leq 3, ~\dim_{R(J)} K^{(i)}(\M)=d\}.$
\end{theorem}
\begin{proof}
First, we consider the following exact sequence induced from \eqref{equation:step2},
	$$0\to [\ker\phi]_1\to {[(R(\M_J)/IR(\M_J))}^\mu]_0\to [R(\M)/(R(\M)_+(1)+R(\M_J))]_1\to 0$$
	and tensor it with $\F_I(J)$ to get the following commutative diagram:
	{\tiny\[
	\xymatrixrowsep{20mm}\xymatrixcolsep{3mm}
	\xymatrix{
		& ([\ker \phi]_1\otimes \F_I(J))(-1) \ar[r] \ar[d]_{\psi_1}  & ({[(R(\M_J)/IR(\M_J))}^\mu]_0\otimes \F_I(J))(-1) \ar[r] \ar[d]^{\psi_2} & ([R(\M)/(R(\M)_+(1)+R(\M_J))]_1 \otimes \F_I(J))(-1) \ar[r]\ar[d]^{\psi_3} & 0 \\
		0 \ar[r] & \ker \phi\ar[r]^{f}  &  [{(R(\M_J)/IR(\M_J)(-1))}^\mu] \ar[r]  & \im{\phi} \ar[r] & 0 
	}
	\]
    }
	of graded $R(J)$-modules where $\psi_2$ is bijective. So, we have an exact sequence \begin{equation}\label{equation:step3}
		0\to K^{(1)}(\M)\to ([R(\M)/(R(\M)_+(1)+R(\M_J))]_1\otimes\F_I(J))(-1)\xrightarrow{\psi_3}\im  \phi\to 0
	\end{equation}
	of graded $R(J)$-modules where $K^{(1)}(\M)=\ker\psi_3$. We put $K^{(3)}(\M)=\L(\M)$ for convenience of notation. Note that $K^{(1)}(\M)$ is a submodule of $ ([R(\M)/(R(\M)_+(1)+R(\M_J))]_1\otimes \F_I(J))(-1)$, $K^{(2)}(\M)$ is a submodule of $ \F_{I_1}(\M_J)$ or $\F_{\m}(\M_J)$  and $K^{(3)}(\M)$ is a submodule of $ S_J(\M)$. So, $\dim K^{(i)}(\M)\leq d$ for $i=1,2,3$.   

By Lemma \ref{lemma-2.2(a)},   $\F_{I_1}(\M_J)\cong M/I_1M[X_1,\ldots,X_d]$ and $\F_{\m}(\M_J)\cong M/\m M[X_1,\ldots,X_d].$ Therefore,  $\ell_A([\F_{I_1}(\M_J)]_n)$ and $\ell_A([\F_{\m}(\M_J)]_n)$ are equal to  $\ell_A(M/M_1)\binom{n+d-1}{d-1}$ for all $n\geq 0.$ 

    
	From \eqref{equation:step1}, \eqref{equation:step1(a)}, \eqref{equation:step2} and \eqref{equation:step3}, we have that 
	\begin{align}
		\ell_A(G(\M)_n)&=\ell_A\Big(\frac{M}{M_1}\Big)\binom{n+d-1}{d-1}+\ell_A\Big(\Big[\frac{R(\M)}{R(\M)_+(1)+R(\M_J)}\Big]_n\Big)-\ell_A([K^{(2)}(\M)]_n)\nonumber\\
		&= \ell_A\Big(\frac{M}{M_1}\Big)\binom{n+d-1}{d-1}+\{\ell_A([\im \phi]_n)+\ell_A\Big(\Big[\frac{S_J(\M)}{\L(\M)}\Big]_{n-1}\Big)\}-\ell_A([K^{(2)}(\M)]_n)\nonumber\\
		&=\ell_A\Big(\frac{M}{M_1}\Big)\binom{n+d-1}{d-1}+\{ \ell_A\Big(\Big[\Big[\frac{R(\M)}{R(\M)_+(1)+R(\M_J)}\Big]_1\otimes \F_I(J)\Big]_{n-1}\Big)-\nonumber\\
		&\qquad\ell_A([K^{(1)}(\M)]_n)\} +\ell_A(S_J(\M)_{n-1})-\ell_A([K^{(3)}(\M)]_{n-1})-\ell_A([K^{(2)}(\M)]_n)\label{equation:step4}
	\end{align}
	for all $n\geq 0$. Since  $ \F_I(J) = R(J)/IR(J) \cong A/I[X_1,\ldots,X_d]$ and 
	$[R(\M)/(R(\M)_+(1)+R(\M_J))]_1=M_1/(M_2+JM)$, we have that 
	\begin{align*}
		\ell_A\Big(\Big[\Big[\frac{R(\M)}{R(\M)_+(1)+R(\M_J)}\Big]_1\otimes \F_I(J)\Big]_{n-1}\Big)&= \ell_A\Big(\frac{M_1}{M_2+JM}\Big)\binom{n-1+d-1}{d-1}\\
		&= \ell_A\Big(\frac{M_1}{M_2+JM}\Big)
		\Big(\binom{n+d-1}{d-1}-\binom{n+d-2}{d-2}\Big).		
	\end{align*}
    Thus,
	\begin{align*}
		\ell_A\Big(\frac{M_n}{M_{n+1}}\Big)&=\ell_A\Big(\frac{M}{M_1}\Big)\binom{n+d-1}{d-1}+\ell_A\Big(\frac{M_1}{M_2+JM}\Big)\binom{n+d-2}{d-1}-\ell_A([K^{(1)}(\M)]_n)\nonumber\\
		&\qquad + \ell_A(S_J(\M)_{n-1})-\ell_A([K^{(3)}(\M)]_{n-1})-\ell_A([K^{(2)}(\M)]_n)\\
		&= \ell_A\Big(\frac{M}{M_1}\Big)\binom{n+d-1}{d-1}+\ell_A\Big(\frac{M_1}{M_2+JM}\Big)\Big\{\binom{n+d-1}{d-1}-\binom{n+d-2}{d-2}\Big\}\\
		&\qquad + \ell_A(S_J(\M)_{n-1})-\ell_A([K^{(1)}(\M)]_n)-\ell_A([K^{(3)}(\M)]_{n-1})-\ell_A([K^{(2)}(\M)]_n)\\
		&=\Big\{ \ell_A\Big(\frac{M}{M_1}\Big)+ \ell_A\Big(\frac{M_1}{M_2+JM}\Big)\Big\}
        \binom{n+d-1}{d-1}-\ell_A\Big(\frac{M_1}{M_2+JM}\Big)\binom{n+d-2}{d-2}\\
		&\qquad +\ell_A(S_J(\M)_{n-1})-\ell_A([K^{(1)}(\M)]_n)-\ell_A([K^{(3)}(\M)]_{n-1})-\ell_A([K^{(2)}(\M)]_n)
	\end{align*}
	for all $n\geq 0$. On comparing the coefficients, we get that
	\begin{align*} 
		e_0(\M)&=\begin{cases}
			\ell_A\Big(\frac{M}{M_1}\Big)+ \ell_A\Big(\frac{M_1}{M_2+JM}\Big)+e_0(S_J(\M))-\mathop\sum\limits_{i\in\Gamma} e_0(K^{(i)}(\M)) & \text{  if } \dim S_J(\M)=d\\        
			\ell_A\Big(\frac{M}{M_1}\Big)+ \ell_A\Big(\frac{M_1}{M_2+JM}\Big)-\mathop\sum\limits_{i\in\Gamma} e_0(K^{(i)}(\M)) & \text{  if } \dim S_J(\M)\leq d-1
		\end{cases}\\
		&= 2\ell_A\Big(\frac{M}{M_1}\Big)+ \ell_A\Big(\frac{M_1}{M_2+JM}\Big)+e_1(\mathcal{M})-e_1^J(M)-e_0(\M)-\mathop\sum\limits_{i\in\Gamma}e_0(K^{(i)}(\M))
	\end{align*}
	since $e_1(\M)-e_0(\M)-e_1^J(M)+\ell(M/M_1)=e_0(S_J(\M))  \text{ if } \dim S_J(\M)=d$ and
	$e_1(\M)-e_0(\M)-e_1^J(M)+\ell(M/M_1)=0  \text{ if } \dim S_J(\M)\leq d-1$
	using Corollary \ref{corr-2.5}. 
\end{proof}

\section*{Author context (source 799--898)}
In view of Lemma \ref{lemma-2.2(a)}, it is easy to observe that $\ass_{R(J)} K^{(i)}(\M)\subseteq \{\m R(J)\}$ for $1\leq i\leq 3.$ To see this, note that 
$K^{(1)}(\M)$ is embedded in $[R(\M)/(R(\M)_+(1)+R(\M_J))]_1\otimes \F_I(J) \cong (M_1/(M_2+JM))\otimes (A/I)[X_1,\ldots,X_d]$ which is a
maximal Cohen-Macaulay $\F_I(J)\cong A/I[X_1,\ldots,X_d]$-module, 
$K^{(2)}(\M)$ is embedded in $ \F^\prime\cong M/M_1[X_1,\ldots,X_d]$
and  $K^{(3)}(\M)$ is embedded in $ S_J(\M)$. Therefore, if $K^{(i)}(\M)\neq 0$, then $\dim K^{(i)}(\M)=d$ for 
$1\leq i\leq 3.$ We will need the following proposition later. Here note that $(C_1)\Rightarrow (C_0)$ and $(C_1)+(C_3)\Rightarrow R(\M_J)$ satisfies $(S_2)$ see \cite[Theorem 6.2]{trung-gcm}. 
\begin{proposition}\label{proposition-3.5}
 Let $M$ be a finitely generated faithful flat $A$-module. Suppose $\M=\M_L$ and $J\subsetneq \m $ or $\M=\M_\I$ and $J\subsetneq I_1$. Suppose the conditions $(C_1)$ and $(C_2)$ are satisfied by $a_1,\ldots,a_d$ for both $A$ and $M$. Also, assume that $(C_3)$ is satisfied by $A$ and $M.$  Then the following conditions are equivalent:
	\begin{enumerate}[label=(\roman*)]
		\item\label{proposition-3.5-equivalent-condition-1}  We have
		$e_1(\mathcal{M})-e_1^J(M)=2e_0(\M)-2\ell_A(M/M_1)-\ell_A(M_1/(M_2+JM)).$
		\item\label{proposition-3.5-equivalent-condition-2} There exist  exact sequences
		$$0\to ((M_1/(M_2+JM))\otimes\F_I(J))(-1)\to R(\M)/(R(\M)_+(1)+R(\M_J))\to S_J(\M)(-1)\to 0,$$
		
		$$0\to \F_{I_1}(\M_J)
        \to G(\M) \to R(\M)/(R(\M)_+(1)+R(\M_J))\to 0$$
         and 
         $$0\to \F_{\m}(\M_J)
        \to G(\M) \to R(\M)/(R(\M)_+(1)+R(\M_J))\to 0$$
		of graded $R(J)$-modules.
		\item \label{proposition-3.5-equivalent-condition-3} $K^{(i)}(\M)=0$ for $1\leq i\leq 3.$
	\end{enumerate}
	When this is the case, we have the injective maps 
	$$\h_{\xi}^i(G(\M))\hookrightarrow \h_{\xi}^i(R(\M)/(R(\M)_+(1)+R(\M_J)))\hookrightarrow \h_{\xi}^i(S_J(\M))(-1)$$
	of graded $R(J)$-modules for $0\leq i\leq d-1$ and 
	$\h_\xi^i(G(\M))\cong \h_{\xi}^i(R(\M)/(R(\M)_+(1)+R(\M_J)))\cong \h_{\xi}^i(S_J(\M))(-1)$ for $0\leq i\leq d-2$.
\end{proposition}
\begin{proof} We show that \ref{proposition-3.5-equivalent-condition-1}$\Rightarrow$ \ref{proposition-3.5-equivalent-condition-3} $\Rightarrow$ \ref{proposition-3.5-equivalent-condition-2}$\Rightarrow$ \ref{proposition-3.5-equivalent-condition-1}.
	\item \ref{proposition-3.5-equivalent-condition-1}  $\Rightarrow$ \ref{proposition-3.5-equivalent-condition-3} 
	Suppose \ref{proposition-3.5-equivalent-condition-1} holds and $K^{(i)}(\M)\neq 0$ for some $1\leq i\leq 3$. Then $i\in\Gamma$ and $e_0(K^{(i)}(\M))\neq 0$ using Theorem \ref{lemma-3.3} which is  a contradiction. 
	
	\item \ref{proposition-3.5-equivalent-condition-3}  $\Rightarrow$ \ref{proposition-3.5-equivalent-condition-2}
	From the exact sequence  \eqref{equation:step3}, we get that $\im \phi\cong ((M_1/(M_2+JM))\otimes\F_I(J))(-1)$. The first exact sequence follows from the exact sequence in \eqref{equation:step2}. Second and third exact sequences follow from the exact sequences in \eqref{equation:step1} and \eqref{equation:step1(a)}, respectively. 
	\item \ref{proposition-3.5-equivalent-condition-2}  $\Rightarrow$ \ref{proposition-3.5-equivalent-condition-1} Suppose \ref{proposition-3.5-equivalent-condition-2}
	holds. Then, for all $n\geq 0$,
	\begin{align*}
		\ell_A(M_n/M_{n+1})=\ell_A(G(\mathcal{M})_n)&=\ell_A(M/M_1)\binom{n+d-1}{d-1}+\ell_A\Big(\Big[\frac{R(\M)}{R(\M)_{+}(1)+R(\M_J)}\Big]_n\Big)\\
		&=\ell_A(M/M_1)\binom{n+d-1}{d-1}+\ell_A\Big(\Big[\frac{M_1}{M_2+JM}\otimes\F_I(J)\Big]_{n-1}\Big)\\
        &\qquad+\ell_A((S_J(\M))_{n-1}).
	\end{align*}
	We get the desired equality using the same argument as done in Theorem \ref{lemma-3.3}.
	
	Finally, since $\F_I(J) \cong (A/I)[X_1,\ldots,X_d]$ by Lemma \ref{lemma-2.2(a)} and  $(M_1/(M_2+JM))\otimes\F_I(J)$ is a maximal 
	Cohen-Macaulay module, we get the injective maps 
	$$\h_\xi^i(G(\M))\hookrightarrow \h_\xi^i(R(\M)/(R(\M)_+(1)+R(\M_J)))\hookrightarrow \h_\xi^i(S_J(\M))(-1)$$
	of graded $R(J)$-modules for $0\leq i\leq d-1$ 
	and $\h_\xi^i(G(\M))\cong \h_{\xi}^i(R(\M)/(R(\M)_+(1)+R(\M_J)))\cong \h_{\xi}^i(S_J(\M))(-1)$ for $0\leq i\leq d-2$.
\end{proof}
\begin{corollary}\label{corollary-3.6} Let $M$ be a finitely generated faithful flat $A$-module and $\M=\M_L$ or $\M=\M_\I.$ Suppose the conditions $(C_1)$ and $(C_2)$  are satisfied by $a_1,\ldots,a_d$ for both $A$ and $M.$ Also, assume that $C_3$ is satisfied by $A$ and $M$ and  $e_1(\M)-e_1^J(M)=2e_0(\M)-2\ell(M/M_1)-\ell(M_1/(M_2+JM))$ holds. Then the following conditions hold:
	\begin{enumerate}[label=(\roman*)]
		\item $M_{n+2}\subseteq J^n M_1$ for all $n\geq 0, $
		\item $J^nM\cap M_{n+1}=J^n M_1$ for all $n\geq 0,$
		\item $\depth G(\mathcal{\M})>0$ and 
		\item $(a_1,\ldots,\check{a_i},\ldots,a_d)M:_Ma_i\subseteq M_2+JM$ for all $1\leq i\leq d.$
	\end{enumerate}
\end{corollary}
\begin{proof}
	If either $J=I_1$ or $J=\m$, then by Remark \ref{remark-I1=Q-case}, $M$ is Cohen-Macaulay $A$-module and $\M=\M_J$. Then the conclusion holds. 
	Suppose $\M=\M_L$ and $J\subsetneq \m $ or $\M=\M_\I$ and $J\subsetneq I_1.$ 
	\begin{enumerate}[label=(\roman*)]
		\item By Proposition \ref{proposition-3.5}, $K^{(3)}(\M)=0$ which gives $M_{n+2}\subseteq J^n M_1 $ for all $n\geq 0$. 
		\item By Proposition \ref{proposition-3.5}, $K^{(2)}(\M)=0$, so $[K^{(2)}(\M)]_n=(J^nM\cap M_{n+1})/J^n M_1=0$ for all $n\geq 0.$
		\item Using the two exact sequences of Proposition \ref{proposition-3.5}\ref{proposition-3.5-equivalent-condition-2}  and depth lemma, it is enough to see that $(M_1/(M_2+JM))\otimes \F_I(J)\cong (M_1/(M_2+JM))\otimes A/I[X_1,\ldots,X_d]$ and $S_J(\M)$ have positive depths as $R(J)$-modules. 
		\item Let $x\in (a_1,\ldots,\check{a_i},\ldots,a_d)M:_Ma_i$. Since $a_1,\ldots,a_d$ satisfies $(C_2)$ for $M$, hence $x\in IM$. Then $xa_i=\mathop\sum\limits_{1\leq j\leq d, j\neq i} a_jx_j$, $x_j\in A$. Here, $x,x_j\in IM\subseteq M_1$ due to the condition $(C_2)$. Let 
		$$g=\overline{x}\otimes \overline{a_it}-\mathop\sum\limits_{1\leq j\leq d, j\neq i} \overline{x_j}\otimes \overline{a_jt}\in (M_1/(M_2+JM))\otimes \F_I(J)$$
		where $\overline{x}, \overline{x_j}$ denote the images of $x,x_j$ in $M_1/(M_2+JM)$ and $\overline{a_it}$ denote the image of $a_it$ in $\F_I(J)$. The image of $g$ under the map $0\to ((M_1/(M_2+JM))\otimes \F_I(J))(-1)\to R(\M)/(R(\M)_+(1)+R(\M_J))$ is 
		$$\overline{xa_it^2}-\mathop\sum\limits_{1\leq j\leq d,j\neq i} \overline{x_ja_jt^2}=0$$
		which implies that $g=0\in (M_1/(M_2+JM))\otimes \F_I(J) \cong (M_1/(M_2+JM))\otimes (A/I)[X_1,\ldots,X_d]$. Therefore, $x\in M_2+JM.$\qedhere
	\end{enumerate}
\end{proof}
\section{\texorpdfstring
{A necessary and sufficient condition for equality in \eqref{very-1st-eqn(a)}}
{A necessary and sufficient condition for equality in the inequality}
}\label{Section-proof-1}
% \section{ A necessary and sufficient condition for equality in    \eqref{very-1st-eqn(a)}}\label{Section-proof-1}
In this section, we prove a necessary and sufficient condition for the equality
$e_1(\M)-e_1^J(M)=2e_0(\M)-2\ell_A(M/M_1)-\ell_A(M_1/(M_2+JM))$ to hold. 
% in Theorem \ref{the-theorem}. 
We begin with the following remark. In what follows, we use the inequality $a(G(\M))+d\leq r_J(\M)$,
whose proof is similar to that of \cite[Proposition 3.2]{nvt} and so we skip providing it here. Recall that 
 $$a(G(\M)):=\max\{n\in \mathbb{Z}~|~ [\h_\xi^d(G(\M))]_n\neq (0)\}.$$

\begin{remark}\label{remark-5.1}Let $M$ be a finitely generated faithful flat $A$-module and $\M= \M_L$ or $\M=\M_\I.$
	Suppose $M_2\subseteq JM$ and  the equality $e_1(\M)-e_1^J(M)=2e_0(\M)-2\ell(M/M_1)-\ell(M_1/(M_2+JM))$ holds. Then
	$M$ is Cohen-Macaulay $A$-module with 
	$M_{n+1}=J^nM_1$ for $n\geq 1$,  $S_J(\M)=(0)$ and $G(\mathcal{M})$ is Cohen-Macaulay $A$-module with $a(G(\M))\leq 1-d.$ 
\end{remark}
\begin{proof}
	By Proposition \ref{proposition:p1},  $e_1(\M)-e_1^J(M)\geq e_0(\M)-\ell(M/M_1)$. Hence,
	\begin{eqnarray*}
		0&=& e_1(\M)-e_1^J(M)-2e_0(\M)+2\ell(M/M_1)+\ell(M_1/JM)\\
		&=&e_1(\M)-e_1^J(M)-2e_0(\M)+\ell(M/M_1)+\ell(M/JM)\\
		&\geq& \ell(M/JM)-e_0(\M)\\
		&\geq& 0.
	\end{eqnarray*}
	Thus, $\ell_A(M/JM)=e_0^J(M)$ which implies that $M$ is Cohen-Macaulay $A$-module, see \cite[Corollary 4.7.11]{bruns-herzog} and $e_1(\M)-e_1^J(M)=e_0(\M)-\ell(M/M_1)$. By Lemma \ref{lemma-2.3(a)} and Lemma \ref{lemma-2.3(b)}, either $S_J(\M)=(0)$ or $\dim S_J(\M)=d$. Using Corollary \ref{corr-2.5}, we get that $S_J(\M)=(0)$.
	 Consequently, $M_{n+1}=J^nM_1$ for
	all $n\geq 1$ and 
	$G(\M)$ is Cohen-Macaulay $A$-module by Vallabrega-Valla criteria. Since $G(\M)$ is Cohen-Macaulay $A$-module, we have $a(G(\M))+d\leq 1$. 
\end{proof}
\section*{Selected author statement}
\begin{proposition}\label{theorem-4.3}
 Let $M$ be a finitely generated faithful flat $A$-module and $\M=\M_L$ or $\M=\M_\I$. Suppose $(C_1)$ and $(C_2)$ conditions are satisfied by $a_1,\ldots,a_d$ for both $A$ and $M$.  Let the equality
	$e_1(\M)-e_1^J(M)=2e_0(\M)-2\ell_A(M/M_1)-\ell_A(M_1/(M_2+JM))$ hold. Then $M_{n+2}\subseteq J^nM_2+W$ for $n\geq 1$ and $W\subseteq M_2$.
\end{proposition}
\section*{Complete author proof}
\begin{proof}
 We first show that $W\subseteq M_2$.  By Lemma \ref{lemma-reduction-steps}(iii), we have that $W\subseteq M_2+JM$ and the equality \eqref{very-1st-eqn(a)} holds in $\overline{M}$. Then by Corollary \ref{corollary-3.6}, $J^n\overline{M}\cap \overline{M_{n+1}}=J^n\overline{M_1}$ for $n\geq 0$. 
 Now let $w=i+q\in W$ where $i\in M_2$ and $q\in JM$. Then $q=w-i\in (JM+W)\cap (M_2+W)=JM_1+W$.
	So, $q\in (JM_1+W)\cap JM=JM_1$ since $JM\cap W=0$ which implies $w=i+q\in M_2$. Hence $W\subseteq M_2$.
	
	For the rest of the proof, we may pass to the ring $\overline{A}$, module $\overline{M}$ and assume that the conditions $(C_1)$ and $(C_2)$ are satisfied by the images of $a_1,\ldots,a_d$ in $\overline{A}$ for both $\overline{A}$ and $\overline{M}.$ Moreover,  $\overline{A}$ and $\overline{M}$ both satisfy $(C_3)$. 
	We show, by induction on $d$, that $M_{n+2}= J^n M_2$ for all $n\geq 1$. Suppose the equality  $e_1(\M)-e_1^J(M)=2e_0(\M)-2\ell_A(M/M_1)-\ell_A(M_1/(M_2+JM))$
	holds. Then
	by Corollary \ref{corollary-3.6}, $\depth G(\M)>0$ and $M_{n+2}\subseteq J^n M_1$ for all $n\geq 0.$
	% % %  {\color{blue}{ We need $Q+W\varsubsetneq I_1+W $ in the statement of the theorem.}} 
	We may assume that 
	$f_1=a_1t\in R(J)$ is a regular element of $G(\M)$.
	
	Suppose $d=1$. Since $M_{n+2}\subseteq J^n M_1\subseteq a_1^nt^nM$ and $f_1^n$ is a regular element of $G(\M)$,  we get that $M_{n+2}=M_{n+2}\cap a_1^nt^nM=a_1^nt^nM_2= J^nM_2$ for all $n\geq 1$.
	
	Suppose $d\geq 2.$ As earlier, let $M^\prime=M/a_1M,$ 
	$\M^\prime=\{M_n^\prime=(M_n+a_1M)/a_1M\}_{n\geq0}$ and $\M_J^\prime=\{(J^nM+a_1M)/a_1M\}_{n\geq 0}$ be the filtration of $M^\prime$. Then the conditions  $(C_1)$ and $(C_2)$ are satisfied for $M^\prime$ and
	$e_1(\M^\prime)-e_1^J(M^\prime)=2e_0(\M^\prime)-2\ell_A(M^\prime/{M_1}^\prime)-\ell_A({M_1}^\prime/({M_2}^\prime+JM^\prime))$ by Lemma \ref{lemma-reduction-steps}.
	Further, $(C_1)$, $(C_2)$ and $(C_3)$ are satisfied for $\overline{M^\prime}=M^\prime/\h_{\m}^0(M^\prime)$. 
	Let $\overline{\M^\prime} =\{(M_n^\prime+\h_{\m}^0(M^\prime))/\h_{\m}^0(M^\prime)\}$ be the $I^\prime \overline{A^\prime}$-good filtration of $\overline{M^\prime}$, where $\overline{A^\prime}=A^\prime/\h_{\m}^0(A^\prime)$. 
	By induction hypothesis, we have for all $n\geq 1$,
	\begin{eqnarray*}
		&M_{n+2}^\prime&\subseteq  J^n M_2 ^\prime+\h_{\m}^0(M^\prime)\\
		\Rightarrow & M_{n+2} & \subseteq J^n M_2+a_1M+\h_{\m}^0(M^\prime)\\
		&& \subseteq J^n M_2+(a_1M:J)
	\end{eqnarray*}
	since $M_{n+2}\subseteq J^nM_1$ by Corollary \ref{corollary-3.6}, $a_1,\ldots,a_d$ is a d-sequence of $M$ and $\h_{\m}^0(M^\prime)=(a_1M:J)$. Now we show that $M_{n+2}\subseteq J^nM_2$ by induction on $n$. let $x\in M_{n+2}$, write $x=y+z$ where $y\in J^nM_2$ and 
	$z\in (a_1M:J)$. Then $z=x-y\in (a_1M:J)\cap M_{n+2}\subseteq (a_1M:J)\cap J^n M_1\subseteq (a_1M:J)\cap JM =a_1M$ since
	$M_{n+2}\subseteq J^nM_1$ by Corollary \ref{corollary-3.6} and $a_1,\ldots,a_d$ is a d-sequence of $M$.  This implies $$z\in a_1M\cap M_{n+2}=a_1M_{n+1}\subseteq JM_{n+1}$$ as $a_1t$ is a regular element for $G(\mathcal{M})$. 
	For $n=1$, we get that $x=y+z\in JM_2$. Suppose $n>1$. Then $z\in JM_{n+1}\subseteq J\cdot J^{n-1}M_2$ by induction hypothesis which implies $x\in J^nM_2$. This completes the proof.
\end{proof}
\par\medskip\noindent\textit{Author citation locators.} Original citation keys and locators are shown. The bibliography database referenced by this fixed author source is not among this unit\textquotesingle s frozen local source inputs; no bibliographic metadata is invented.
\end{document}
